\documentclass[12pt]{article}

\usepackage{fullpage}
\usepackage{multicol,multirow}
\usepackage{tabularx}
\usepackage{ulem}
\usepackage[utf8]{inputenc}
\usepackage{mathtools}
\usepackage{pgfplots}
\usepackage[russian]{babel}

\begin{document}

\section*{Лабораторная работа №\,4 по курсу дискрeтного анализа: строковые алгоритмы}

Выполнил студент группы М8О-208Б-22 МАИ \textit{Немкова Анастасия}.

\subsection*{Условие}
Необходимо реализовать один из стандартных алгоритмов поиска образцов для указанного алфавита.\newline
\newline
\textbf{Вариант алгоритма:}
Поиск одного образца-маски: в образце может встречаться «джокер» (представляется символом ? — знак вопроса), равный любому другому символу. При реализации следует разбить образец на несколько, не содержащих «джокеров», найти все вхождения при помощи алгоритма Ахо-Корасик и проверить их относительное месторасположение.\newline
\newline
\textbf{Вариант алфавита:}
Слова не более 16 знаков латинского алфавита (регистронезависимые)


\subsection*{Метод решения}

Алгоритм Ахо-корасик - это алгоритм поиска подстроки в строке, который реализует поиск множества подстрок из словаря в данной строке.

Для начала из всех подстрок образца, не содержащих джокеров, строим trie - дерево, содержащее корень, равный пустой строке, рёбра, по которым можно пеpeйти к следующей вершине по символу, и другие вершины, являющиеся префиксами образца. 

Затем для каждой вершины данного дерева строим связи неудач, обходя дерево в ширину. Связь неудачи связывает две вершины дерева - первая, которая получена в результате прохода по тексту и совпадения символов текста с образцом до вершины, в которой найдено несовпадение; вторая - вершина другого образца, префикс которого является максимальным суффиксом пройденного пути. Для корня и вершин первого уровня суффиксной ссылкой будет ссылка на корень. Для любой другой вершины
нужно перейти по суффиксной ссылке её родителя в узел $n$. Если из узла $n$ можно перейти в другой узел $x$ по последнему символу
строки исходной вершины, тогда ставим суффиксную ссылку на узел $x$. Если нельзя перейти, тогда переходим по суффиксной ссылке
узла $n$, пока не упрёмся в корень и делаем то же самое. Если дошли до корня, тогда ставим суффиксную ссылку на корень. Также при построении дерева нужно найти длины подстрок образца, не содержащих джокеров, и их начальные позиции в образце для дальнейшей проверки местоположения подстрок.

При поиске необходимо пройти по тексту и найти вхождения каждого подобразца в текст. Для каждого такого начального вхождения мы увеличиваем счетчик в ячейке [текущая позиция в тексте - позиция подобразца в образце] вектора, равного длине текста и инициализированного нулями, на единицу. В итоге образец считается найденным в тексте, если на i-ой позиции данного вектора находится значение, равное количеству подобразцов без джокеров.


\subsection*{Описание программы}

Для реализации словаря был создан класс дерева Trie, каждый узел которого содержит ассоциативный массив для хранения дочерних узлов текущего узла, Указатель на узел, к которому происходит переход в случае неудачного совпадения текущего символа, и вектор индексов подобразцов, которые заканчиваются в данном узле. Само дерево хранит указатель на корень, вектор с длинами подобразцов и счетчик подобразцов без джокеров. В классе Trie реализованы следующие методы:

\begin{itemize}
    \item void Create(const std::vector<std::string>\& patterns) - построение дерева на основе входного образца.
    \item void Search(const std::vector<std::string>\& text, const int\& patSize, std::vector<std::pair<int, int>>\& answer) - поиск подобразцов в тексте.
    \item void CreateLinks() - создание ссылок на неудачные переходы.
\end{itemize}

\subsection*{Дневник отладки}

\begin{enumerate}
    \item 23 май 2024, 02:16:59 OK на 1 тесте
\end{enumerate}

\subsection*{Тест производительности}

Для измерения производительсти сравнивается время работы реализованного алгоритма Ахо-Корасик и наивного поиска с помощью стандартного контейнера std::find. Подсчет времени производился с помощью библиотеки chrono, которая позволяет фиксировать время в начале и конце выполнения сортировки. Первый тест состоит из $10^2$ строк по 50 слов в каждой и 20 слов в паттрене, второй - из $10^3$ строк по 500 слов в каждой и 100 слов в паттрене и третий из $10^4$ строк по 1000 слов в каждой и 100 слов в паттрене.

На графике представлена зависимость времени выполнения поиска от объёма входных данных.


\begin{figure}[htbp]
    \centering
    \begin{tikzpicture}
        \begin{axis}[
            xlabel={Объём входных данных ($\times 10^3$)},
            ylabel={Время работы (sec)},
            grid=major,
            xmin=0, xmax=10,
            ymin=0, ymax=3,
            xtick={0, 1, 2, 3, 4, 5, 6, 7, 8, 9, 10},
            ytick={0, 0.3, 0.6, 0.9, 1.2, 1.5, 1.8, 2.1, 2.4, 2.7, 3},
            legend style={at={(0.5,-0.2)},anchor=north},
            legend columns=2,
            width=0.8\textwidth,
            height=0.5\textwidth,
            ]
            \addplot[color=blue,mark=*] coordinates {
                (0,0)
                (0.1, 0.00036)
                (1, 0.0321)
                (10, 0.6513)
            };
            \addlegendentry{Ахо-Корасик}
            \addplot[color=red,mark=square] coordinates {
                (0,0)
                (0.1, 0.00034)
                (1, 0.1074)
                (10, 2.0325)
            };
            \addlegendentry{std::find}
        \end{axis}
    \end{tikzpicture}
    \label{fig:graph}
\end{figure}

\newpage
\subsection*{Выводы}

В ходе данной лабораторной работы был изучен алгоритм Ахо-Корасика, позволяющий искать несколько образцов в строке за линейное время. Также была создана структура дерева Trie и реализован автомат для поиска по этому дереву. 

Время требуемое на построение дерева — O(n), где n — суммарная длина всех подобразцов, не содержащих джокера. Если число джокеров ограничено фиксированной постоянной, то решение задачи поиска возможно за линейное время O(n + m), где n - суммарная длина всех подстрок образца без джокеров, m — длина текста. Для отображения образца в дерево используется map, реализованая на красно-черном дереве. Время обращения к элементам пропорционально логарифму числа элементов, а это константа. Время поиска вхождений составляет O(m + t), t — число вхождений. Число джокеров ограничено, независимо от размера образца, поэтому алгоритм работает за линейное время.


\end{document}

